\chapter{Coarse Moduli Spaces}

In this chapter we introduce coarse moduli spaces, which provide a way to pass
from algebraic stacks to algebraic spaces by identifying objects up to
isomorphism. This construction plays a central role in the Keel-Mori theorem.

\section{Motivation}

Let $\mathcal{X}$ be an algebraic stack. For a scheme $T$, the groupoid
$\mathcal{X}(T)$ consists of objects over $T$ together with their isomorphisms.
A natural goal is to construct a geometric object whose points correspond to
isomorphism classes of objects of $\mathcal{X}$.

However, in general $\mathcal{X}$ is not representable by a scheme or algebraic space due to the presence of automorphisms. The notion of a coarse moduli space provides a best possible approximation.

\section{Definition of Coarse Moduli Space}

\begin{definition}[Coarse moduli space]
	\label{def:coarse-moduli-space}
	Let $\mathcal{X}$ be an algebraic stack. A morphism
	\[
	\pi : \mathcal{X} \to X
	\]
	to an algebraic space $X$ is called a coarse moduli space if:
	
	\begin{enumerate}
		\item \textbf{Universality:}  
		For every algebraic space $Y$ and every morphism
		\[
		f : \mathcal{X} \to Y,
		\]
		there exists a unique morphism
		\[
		g : X \to Y
		\]
		such that $f = g \circ \pi$.
		
		\item \textbf{Geometric points:}  
		For every algebraically closed field $k$, the induced map
		\[
		|\mathcal{X}(k)| \to X(k)
		\]
		identifies $X(k)$ with the set of isomorphism classes of objects of
		$\mathcal{X}(k)$.
	\end{enumerate}
\end{definition}

\begin{remark}
	The coarse moduli space forgets all automorphism data and only remembers
	isomorphism classes. In particular, it does not carry a universal family in
	general.
\end{remark}


\section{Basic Properties}



\section{Coarse Moduli as Quotients}

\begin{definition}[Coarse quotient]
	(\cite[][Tag 04A1]{stacks-project})
	Let $R \rightrightarrows U$ be a groupoid in algebraic spaces. A morphism
	\[
	\phi : U \to X
	\]
	is called a \emph{coarse quotient} if:
	\begin{enumerate}
		\item $\phi$ is a categorical quotient,
		\item $\phi$ identifies geometric points with orbits.
	\end{enumerate}
\end{definition}

\begin{remark}
	Coarse moduli spaces can be viewed as coarse quotients of groupoids arising
	from presentations of stacks.
\end{remark}


\section{Examples}

\begin{example}[Trivial case]
	If $\mathcal{X} = X$ is an algebraic space, then the identity
	\[
	X \to X
	\]
	is a oarse moduli space.
\end{example}

\begin{example}[Quotient by a finite algebraic group]
	Let $G$ be a finite group acting on a scheme $U$. Then the quotient stack
	\[
	[U/G]
	\]
	has a coarse moduli space given by the geometric quotient
	\[
	U \to U/G.
	\]
\end{example}

\section{Existence Problem}

\begin{remark}
	In general, coarse moduli spaces do not always exist. However, the
	Keel-Mori theorem shows that they exist for a large class of algebraic stacks,
	namely those with finite inertia.
\end{remark}
